% SGA 3, 第 VI_B 篇, 第 1 轮英文重构.
% 范围: 引理 5.6.2.0 全文, 含底本编者注 42.
% 底本: Polo--Gille Exp6B-13oct24.pdf, 分篇第 31--32 页
% (合订本第 372--373 页; 原版页边印刷页码 354, 重版本页眉 361--362).
% 推论 5.6.2 始于分篇第 32 页 / 合订本第 373 页.
% OCR 与 Reinhold 英文 Markdown 仅为比对材料.
% 本单元无图.

\medskip
\phantomsection\label{lem:sga3-VIB-5.6.2.0}%
\noindent\textbf{引理 5.6.2.0.---}%
\footnote{\label{note:sga3-VIB-42}编者注: 此引理为我们新增,
由 O.\ Gabber 告知. 它使我们能够简化
\SGARefLink{cor:sga3-VIB-5.6.2}{5.6.2} 的证明,
并以更一般的形式证明\hyperref[thm:sga3-VIB-5.3]{定理 5.3}
及 \hyperref[cor:sga3-VIB-5.4]{5.4};
见下面的 \hyperref[thm:sga3-VIB-5.7]{5.7} 和 \hyperref[sec:sga3-VIB-5.8]{5.8}.}
\emph{设 \(k\) 是域, \(G\) 是非空有限型 \(k\)-概形.}
\begin{enumerate}
\item[\textup{\SGAInlineRefTarget{lem:sga3-VIB-5.6.2.0:item:i}(i)}]
  \emph{设 \(K\) 是 \(k\) 的扩张, \(W\) 是 \(G_K\) 的稠密开集.
  用 \(\pi\) 表示投影 \(G_K\to G\). 则}
  \[
    U=\bigl\{g\in G\mathrel{\big|}W\text{ 包含 }\pi^{-1}(g)\text{ 的每个极大点}\bigr\}
  \]
  \emph{是 \(G\) 的稠密开集.
  \textup{(注意: 若 \(g\) 是 \(G\) 的一个属于 \(U\) 的闭点,
  则 \(\pi^{-1}(g)\subset W\).}}

\item[\textup{\SGAInlineRefTarget{lem:sga3-VIB-5.6.2.0:item:ii}(ii)}]
  \emph{再假设 \(G\) 几何不可约. 设}
  \[
    \mu:G\times X\longrightarrow Y
  \]
  \emph{是 \(k\)-概形态射, \(x\) 是 \(X\) 的点,
  \(\Omega\) 是 \(Y\) 的开集, 满足}
  \[
    \mu_x^{-1}\bigl(\Omega_{\kappa(x)}\bigr)\ne\varnothing,
  \]
  \emph{其中 \(\mu_x\) 表示态射}
  \[
    G_{\kappa(x)}\longrightarrow Y_{\kappa(x)},
    \qquad g\longmapsto\mu(g,x).
  \]
  \emph{则}
  \[
    U=\bigl\{g\in G\mathrel{\big|}\mu\text{ 将 }g\times x\text{ 的每个极大点映入 }\Omega\bigr\}
  \]
  \emph{是 \(G\) 的稠密开集.
  对 \(G\) 的每个属于 \(U\) 的闭点 \(g\), 有}
  \[
    \mu(g\times x)\subset\Omega,
  \]
  \emph{因而对每个位于 \(g\) 上的点 \(g'\in G_{\kappa(x)}\)
  \textup{(分别地, 位于 \(x\) 上的点 \(x'\in X_{\kappa(g)}\))}, 有}
  \[
    \mu(g',x)\in\Omega_{\kappa(x)}
    \quad\text{\textup{(分别地, }}\,
    \mu(g,x')\in\Omega_{\kappa(g)}\text{\textup{)}},
  \]

\item[\textup{\SGAInlineRefTarget{lem:sga3-VIB-5.6.2.0:item:iii}(iii)}]
  \emph{假设 \(G=H^0\), 其中 \(H\) 是局部有限型 \(k\)-群概形,
  它作用于非空 \(k\)-概形 \(X\), 使态射}
  \[
    H\times_S X\longrightarrow X\times_S X,
    \qquad (h,x)\longmapsto(hx,x)
  \]
  \emph{满射. 设 \(U\) 是 \(X\) 的开集. 则:}
  \begin{enumerate}
  \item[\textup{\SGAInlineRefTarget{lem:sga3-VIB-5.6.2.0:item:a}(a)}]
    \emph{\(G\cdot U\) 是 \(X\) 的开集,
    等于 \(X\) 中泛点属于 \(U\) 的不可约分支之并.}
  \item[\textup{\SGAInlineRefTarget{lem:sga3-VIB-5.6.2.0:item:b}(b)}]
    \emph{假设 \(U\) 在 \(X\) 的每个不可约分支中稠密.
    则对 \(X\) 的每个有限子集 \(A\), 存在闭点 \(g\in G\), 使得}
    \[
      g\cdot A'\subset U_{\kappa(g)},
    \]
    \emph{其中 \(A'\) 是 \(A\) 在 \(X_{\kappa(g)}\) 中的逆像.
    特别地, 态射}
    \[
      G\times U\longrightarrow X
    \]
    \emph{满射.}
  \end{enumerate}
\end{enumerate}

\emph{证明.}
\textup{\SGARefLink{lem:sga3-VIB-5.6.2.0:item:i}{(i)}}
设 \(\bar{k}\) 是 \(k\) 的可分闭包,
\(L\) 是 \(k\) 的一个同时包含 \(\bar{k}\) 与 \(K\) 的副本的扩张.
用 \(\pi_L\) (分别地, \(\phi\)) 表示投影
\[
  G_L\longrightarrow G
  \qquad\text{(分别地, \(G_L\longrightarrow G_K\)).}
\]
对每个 \(g\in G\), 映射 \(\phi\) 将 \(\pi_L^{-1}(g)\)
的极大点满射地映到 \(\pi^{-1}(g)\) 的极大点. 因此,
\[
  U=\bigl\{g\in G\mathrel{\big|}W_L\text{ 包含 }\pi_L^{-1}(g)\text{ 的每个极大点}\bigr\}.
\]
于是将 \(K\) 改为 \(L\), 即归结为 \(K\) 包含 \(\bar{k}\) 的情形.

由于投影
\[
  p:G_K\longrightarrow G_{\bar{k}}
\]
满射且开, \(V=p(W)\) 是 \(G_{\bar{k}}\) 的稠密开集. 此外,
\[
  p^{-1}(g)=\operatorname{Spec}
    \bigl(\kappa(g)\otimes_{\bar{k}}K\bigr)
\]
对每个 \(g\in G_{\bar{k}}\) 都不可约
(参见 EGA~IV\(_2\), 4.3.3).
所以对每个 \(g\in V\), \(p^{-1}(g)\) 的泛点属于 \(W\).

另一方面, 置
\[
  \mathscr G=\operatorname{Gal}(\bar{k}/k).
\]
投影 \(q:G_{\bar{k}}\to G\) 满射,
且 \(\mathscr G\) 在 \(q\) 的纤维上传递作用. 因此
\[
  U=q(V'),
\]
其中 \(V'\) 是 \(V\) 的各个 \(\mathscr G\)-共轭的交.

设 \(Z\) 是闭集 \(G_{\bar{k}}-V\), 赋予其约化闭子概形结构.
由于 \(G_{\bar{k}}\), 因而 \(Z\), 在 \(\bar{k}\) 上有限呈示,
概形 \(Z\) 由某个 \(k\) 的有限 Galois 扩张 \(k_1\) 上的约化闭子概形
\[
  Z_1\subset G\otimes_k k_1
\]
作基变换得到.
于是 \(Z\) 的 \(\mathscr G\)-共轭只有有限多个,
它们的并仍是 \(G_{\bar{k}}\) 的无处稠密闭集 \(F\). 可知
\[
  V'=G_{\bar{k}}-F
\]
是 \(G_{\bar{k}}\) 的稠密开集.
由于 \(q:G_{\bar{k}}\to G\) 满射且开,
\(U=q(V')\) 因而是 \(G\) 的稠密开集.

此外, 对 \(G\) 的每个闭点 \(g\),
纤维 \(\pi^{-1}(g)\) 由 \(G_K\) 的有限个闭点组成.
所以若 \(g\in U\), 则
\[
  \pi^{-1}(g)\subset W.
\]
这证明了 \textup{\SGARefLink{lem:sga3-VIB-5.6.2.0:item:i}{(i)}},
而 \textup{\SGARefLink{lem:sga3-VIB-5.6.2.0:item:ii}{(ii)}} 由它得出.

断言 \textup{\SGARefLink{lem:sga3-VIB-5.6.2.0:item:iii}{(iii)}\SGARefLink{lem:sga3-VIB-5.6.2.0:item:a}{(a)}}
已在 VI\(_A\), \SGARefLink{prop:sga3-VIA-2.6.4}{2.6.4} 中证明.
最后, 若 \(U\) 在 \(X\) 的每个不可约分支中稠密, 则
\[
  X=G\cdot U.
\]
故对每个 \(x\in X\),
\[
  \mu_x^{-1}\bigl(U_{\kappa(x)}\bigr)
\]
是 \(G_{\kappa(x)}\) 的非空开集,
而 \textup{\SGARefLink{lem:sga3-VIB-5.6.2.0:item:iii}{(iii)}\SGARefLink{lem:sga3-VIB-5.6.2.0:item:b}{(b)}}
由 \textup{\SGARefLink{lem:sga3-VIB-5.6.2.0:item:ii}{(ii)}} 得到.

% 翻译覆盖：本组件全文；源文件 44_expose_VIB_lemma5620_en.tex。
% 自检：行内数学 104，陈列数学 21，原生 tikzcd 0；全部数学及编号保留，公式内自然语言已译。
% 引用键核对：label 2，hyperref 4，SGARefLink 10，行内目标 5；脚注 1，星号脚注正文 0，均与源文件一致。
% 编译及渲染审查由主线程整合完成；此处未宣称编译或独立校对完成。
